% ID: FIS-OND-DOP-002
% Titolo: Raddoppio della frequenza percepita davanti a una sorgente
% Disciplina: Fisica
% Area: Onde
% Argomento: Effetto Doppler
% Sottoargomento: Sorgente sonora in moto
% Classe: Quarta liceo scientifico
% Difficolta: 3
% Tipo: problema_numerico
% Risultato: v_s = v_suono / 2; non dipende dalla frequenza
% Tempo_stimato: 8 min
% Competenze: modellizzazione, ragionamento su rapporti, interpretazione fisica
% Prerequisiti: effetto Doppler, frequenza, velocita del suono
% Tag: fisica, onde, suono, effetto_doppler, sorgente_in_moto
% Fonte: verifica_fisica_onde_anonimizzata
% Autore: Riccardo Carli
% Licenza: CC-BY-SA-4.0

\begin{esercizio}
A che velocita si dovrebbe muovere un aereo per far si che il rombo delle sue
turbine raddoppi in frequenza per gli osservatori che si trovano davanti ad
esso sul suo cammino?

Il risultato dipende dalla frequenza considerata?
\end{esercizio}

\begin{soluzione}
Per una sorgente sonora che si avvicina a osservatori fermi, la frequenza
osservata davanti alla sorgente e:
\[
f'=f\frac{v}{v-v_s},
\]
dove \(v\) e la velocita del suono e \(v_s\) la velocita della sorgente.

Vogliamo che la frequenza raddoppi:
\[
f'=2f.
\]
Quindi:
\[
2f=f\frac{v}{v-v_s}.
\]
Semplificando \(f\):
\[
2=\frac{v}{v-v_s}.
\]

\[
2(v-v_s)=v
\qquad\Rightarrow\qquad
v_s=\frac{v}{2}.
\]

Se si assume \(v=343\,\mathrm{m/s}\), allora:
\[
v_s\simeq 1{,}7\cdot 10^2\,\mathrm{m/s}.
\]

Il risultato non dipende dalla frequenza emessa, perche \(f\) si semplifica nel
rapporto tra frequenza osservata e frequenza emessa.
\end{soluzione}
